En esta sección se explicará qué se ha medido, por qué, y las limitaciones encontradas.

\begin{itemize}
    \item \textbf{SI-SDR:}
Como métrica principal se ha utilizado \textbf{SI-SDR} ya que es una métrica habitual en separación de fuentes y permite comparar la señal estimada
con la señal de referencia. Mide cuánta parte de la estimación se alinea con la fuente objetivo, frente a la parte que puede considerarse distorsión.

A la hora de la interpretación de los resultados, para SI-SDR, se valora a mejor un valor más alto de esta métrica mientras que un valor más bajo implica peor separación.

\[
\mathrm{SI\text{-}SDR} = 10 \log_{10}\frac{\|s_{\mathrm{target}}\|^2}{\|e_{\mathrm{noise}}\|^2}
\]

Donde: 
\begin{itemize}
    \item $s_{\mathrm{target}}$ es la señal de referencia.
    \item $e_{\mathrm{noise}}$ es el error residual, es decir, todo lo que queda en la estimación y que no pertenece a $s_{\mathrm{target}}$.
\end{itemize}

    \item \textbf{Baseline de mezcla:}
Un modelo no se considera útil si no supera la línea base de mezcla en el conjunto de test. Por ello, se ha añadido un \textit{baseline} trivial, donde se utiliza la mezcla original como estimación de cada fuente para así detectar modelos
que generan audio, pero no separan de verdad.
Para ello se ha definido una línea base trivial (\textit{baseline}) que consiste en utilizar la mezcla original como estimación de cada fuente objetivo. Su \textit{SI-SDR} se calcula del mismo modo que el de los modelos.El valor de \textit{SI-SDRi} expresa la diferencia entre el \textit{SI-SDR} del modelo y esta baseline. Un SI-SDRi positivo indica que el modelo mejora respecto a utilizar directamente la mezcla como estimación y uno negativo indica que empeora.

Se han utilizado algunas métricas de agregación para tener información adicional sobre el rendimiento de los procesos realizados, así como:

\begin{itemize}
    \item Media: da una medida de rendimiento promedio para un conjunto de valores de, por ejemplo SI-SDR.
    \item Mediana: da un valor de robustez frente a \textit{outliers}.
    \item Desviación típica: determina la distancia a la que están la mayoría de los valores de la media y se ha utilizado para valorar la estabilidad entre pistas.
\end{itemize}

\item \textbf{Métricas complementarias:}

Se han utilizado también varias métricas complementarias a modo de depuración y de diagnóstico: correlación entre fuentes de salida y mezcla, estadísticas de máscara y matriz de correlación estimación-referencia.

    
\end{itemize}














